\section{The model toolkit}

\begin{frame}{Model families at a glance}
  \footnotesize
  \begin{tabular}{p{0.2\textwidth}p{0.22\textwidth}p{0.27\textwidth}p{0.22\textwidth}}
    \toprule
    Family & What drives the rates & Key result & Checked here \\
    \midrule
    Zero intelligence, CST & distance to opposite quote & depth and spread set by $\lambda/\theta$, $\mu/\lambda$ & depth $\to\lambda/\theta$ far out \\
    Best queues (CdL) & constants & $p_\uparrow(q_b,q_a)$; $1/t$ duration tail; $\sigma^2 \propto \lambda\delta^2$ & MC and exact solve \\
    Queue-reactive (HLR) & queue size & stationary law in product form & TV distance \EthreeTV \\
    Hawkes & past events & $\bar\Lambda = (I-G)^{-1}\Lambda^0$; signature plot & MLE, KS, $C(h)$ \\
    Propagator & past trade signs & diffusive only at $\zeta_c = (1-\gamma_\varepsilon)/2$ & slopes $\pm 0.01$ \\
    Market making & control of quotes & AS and exact GLFT quotes & CE ranking \\
    \bottomrule
  \end{tabular}
  \vspace{0.5em}

  \small CST: \citet{cont2010stochastic}. CdL: \citet{cont2013price}. HLR: \citet{huang2015simulating}. GLFT: \citet{gueant2013dealing}. CE: certainty equivalent. MC: Monte Carlo. TV: total variation. MLE: maximum likelihood. KS: Kolmogorov-Smirnov test.
\end{frame}

% ---------------------------------------------------------------- zero intelligence / CST
\begin{frame}{Zero intelligence and the Markovian book}
  \begin{columns}[T]
    \begin{column}{0.5\textwidth}
      \small
      State: signed queue sizes $X_p$ on a price grid. With $i$ the distance (ticks) to the \emph{opposite} best quote:
      \begin{itemize}
        \item limit orders at distance $i$: rate $\lambda(i)$
        \item market orders: rate $\mu$ per side
        \item each resting order cancels at rate $\theta(i)$
      \end{itemize}
      Zero intelligence \citep{smith2003statistical}: $\lambda,\theta$ constant. CST: $\lambda(i) = \lambda_1 i^{-a}$, with $\lambda_1$ (s$^{-1}$) the rate at distance one and $a$ a dimensionless exponent.

      \vspace{0.4em}
      \textbf{Dimensional analysis.} With prices in ticks and counts in orders, every stationary statistic depends only on
      \[ \lambda/\theta \ \text{(depth)}, \quad \mu/\lambda \ \text{(price scale)}, \quad L \ \text{(window)}. \]
    \end{column}
    \begin{column}{0.48\textwidth}
      \small
      Why it matters: \citet{farmer2005predictive} found that this model, with rates measured per stock, explains most of the cross-sectional variation in mean spread on the London Stock Exchange.

      \vspace{0.4em}
      What it misses: rates depend on the state and on history; order signs have no memory; no information.

      \vspace{0.4em}
      Simulated exactly by the Gillespie algorithm on a moving grid.
    \end{column}
  \end{columns}
\end{frame}

\begin{frame}{Markovian book: depth profile and spread}
  \centering
  \includegraphics[width=0.92\textwidth]{fig1_markov_book.png}
  \small
  \begin{itemize}
    \item Far from the quotes, depth $\bar Q(i)$ (orders) follows the M/M/$\infty$ value $\lambda(i)/\theta(i)$; near the quotes, market orders thin the book.
    \item Mean spread: $\EoneCSTspread$ ticks (power law) against $\EoneZIspread$ ticks (zero intelligence, $\mu/\lambda$ six times larger). Standard errors from 20 batch means over $\EoneNevents$ events.
  \end{itemize}
\end{frame}

% ---------------------------------------------------------------- Cont-de Larrard
\begin{frame}{Best queues: direction and timing of the next price move}
  \small
  Two best queues with one-tick spread; each grows at rate $\lambda$ and shrinks at rate $\nu$ (market orders plus cancellations, s$^{-1}$). The price moves up when the ask queue empties first \citep{cont2013price}.
  \begin{columns}[T]
    \begin{column}{0.52\textwidth}
      First passage of a queue of size $q$ to zero, density (s$^{-1}$):
      \[ f_q(t) = \frac{q}{t}\Big(\frac{\nu}{\lambda}\Big)^{q/2} I_q\big(2\sqrt{\lambda\nu}\,t\big)\,e^{-(\lambda+\nu)t}, \]
      with $I_q$ the modified Bessel function. Then
      \[ p_\uparrow(q_b,q_a) = \int_0^\infty f_{q_a}(t)\,P(\tau_{q_b} > t)\,\dd t . \]
    \end{column}
    \begin{column}{0.46\textwidth}
      Balanced case $\lambda = \nu$, large queues:
      \[ p_\uparrow \approx \tfrac{2}{\pi}\arctan(q_b/q_a). \]
      Time to the next move: $P(\tau > t) \sim q_b q_a/(\pi\lambda t)$, so the \hl{mean duration is infinite}.

      Price diffusion, with $M_F$ the mean product of regenerated queue sizes (orders$^2$):
      \[ \sigma^2 = \pi\lambda\delta^2/M_F \quad (\text{time scale } t\,n\log n). \]
    \end{column}
  \end{columns}
\end{frame}

\begin{frame}{Best queues: closed forms against simulation}
  \centering
  \includegraphics[width=0.92\textwidth]{fig2_best_queues.png}
  \small
  \begin{itemize}
    \item Monte Carlo ($\EtwoNmc$ paths per point) matches the closed form within $\EtwoMaxDev$ standard errors; an exact sparse linear solve agrees to $2\times10^{-3}$.
    \item The arctangent law is off by $\EtwoArctanErrSmall$ at $(2,5)$ and $\EtwoArctanErr$ at $(40,20)$. Imbalance alone is a strong predictor, which is why micro-prices work \citep{stoikov2018micro}.
  \end{itemize}
\end{frame}

% ---------------------------------------------------------------- queue reactive
\begin{frame}{Queue-reactive models}
  \small
  \citet{huang2015simulating} estimate every rate as a function of the queue size $q$ (in average event sizes): $q\to q+1$ at $\lambda^L(q)$, $q\to q-1$ at $\lambda^C(q)+\lambda^M(q)$. Model I (independent queues) has the product-form stationary law
  \[ \pi(q) = \pi(0)\prod_{j=1}^{q}\lambda^L(j-1)\big/\big[\lambda^C(j)+\lambda^M(j)\big]. \]
  \centering
  \includegraphics[width=0.58\textwidth]{fig3_queue_reactive.png}

  \raggedright\footnotesize
  Illustrative rates (a); product formula against one simulated queue of $\EthreeNevents$ events (b): total variation distance $\EthreeTV$. Models II and III couple the best queues and move a reference price; \citet{wu2019queue} add Hawkes terms.
\end{frame}

% ---------------------------------------------------------------- Hawkes
\begin{frame}{Hawkes processes: order flow excites itself}
  \small
  For event types $u, v$ (buy, sell, cancel, \ldots) the intensity (s$^{-1}$) is
  \[ \Lambda_u(t) = \Lambda^0_u + \sum_v \sum_{t^v_j < t} \alpha_{uv}\,e^{-\beta_{uv}(t - t^v_j)}, \]
  with baseline $\Lambda^0_u$, jump $\alpha_{uv}$ and decay $\beta_{uv}$, all in s$^{-1}$. The branching matrix $G_{uv} = \alpha_{uv}/\beta_{uv}$ counts direct offspring; stationarity needs spectral radius $\rho(G) < 1$, and then $\bar\Lambda = (I - G)^{-1}\Lambda^0$.
  \begin{columns}[T]
    \begin{column}{0.5\textwidth}
      \begin{itemize}
        \item Fitted order flow is often \hl{near critical}, $\rho(G)\approx 1$ \citep{hardiman2013critical}, and the kernels decay slowly, so they must be estimated from about 10 $\mu$s to 100 s \citep{bacry2016estimation}.
        \item Exact likelihood in one pass for exponential kernels; residuals $\mathcal{A}_u(t_j) - \mathcal{A}_u(t_{j-1})$ of the compensator $\mathcal{A}_u$ are i.i.d.\ unit exponential under the true model.
      \end{itemize}
    \end{column}
    \begin{column}{0.48\textwidth}
      \begin{itemize}
        \item Nearly unstable Hawkes processes have diffusive scaling limits \citep{jaisson2015limit}; with heavy-tailed kernels the limits are \hl{rough}, which gives rough volatility with leverage from microstructure \citep{eleuch2018microstructural,wang2026rough}.
        \item Estimation pitfalls (seasonality, regime shifts, one-scale kernels): deep dive 4.
      \end{itemize}
    \end{column}
  \end{columns}
\end{frame}

\begin{frame}{Hawkes: fitting buy and sell market orders}
  \centering
  \includegraphics[width=0.8\textwidth]{fig4_hawkes.png}
  \begin{columns}[T]
    \begin{column}{0.48\textwidth}
      \footnotesize
      Bivariate exponential Hawkes model of buys and sells, $T_{\mathrm{obs}} = \EfourT$ s of data: $\EfourNbuy$ buys and $\EfourNsell$ sells, true spectral radius $\rho(G) = \EfourRho$. (a) Intensities over one minute; ticks mark events.
    \end{column}
    \begin{column}{0.48\textwidth}
      \footnotesize
      \begin{itemize}
        \item MLE recovers all ten parameters within $\EfourMaxZ$ standard errors (observed information).
        \item (b) Rescaled waiting times are unit exponential: Kolmogorov-Smirnov $p = \EfourKSbuy$ and $\EfourKSsell$.
      \end{itemize}
    \end{column}
  \end{columns}
\end{frame}

\begin{frame}{Hawkes price model: microstructure noise}
  \begin{columns}[c]
    \begin{column}{0.48\textwidth}
      \small
      Up ticks $N_+$ and down ticks $N_-$ form a symmetric Hawkes process: self jump $\alpha_s$, cross jump $\alpha_c$, decay $\beta$ \citep{bacry2013modelling}. With $g = (\alpha_s - \alpha_c)/\beta$ and $\eta = \beta(1-g)$ (s$^{-1}$), the variance per unit time at sampling interval $h$ (s) is
      \[ C(h) = 2\bar\Lambda\Big[\tfrac{1}{(1-g)^2} + \Big(1 - \tfrac{1}{(1-g)^2}\Big)\tfrac{1 - e^{-\eta h}}{\eta h}\Big] \]
      in ticks$^2$ s$^{-1}$. Cross-excitation ($g<0$) gives the falling signature plot of tick data: bid-ask bounce. Derivation in the backup slides.

      \vspace{0.3em}
      Simulation matches within $\EfiveMaxZ$ standard errors over 42 points ($\EfiveNrep$ replicates each).
    \end{column}
    \begin{column}{0.5\textwidth}
      \centering
      \includegraphics[width=0.95\textwidth]{fig5_signature.png}
    \end{column}
  \end{columns}
\end{frame}

% ---------------------------------------------------------------- propagator / impact
\begin{frame}{The propagator model: memory and efficiency}
  \begin{columns}[T]
    \begin{column}{0.5\textwidth}
      \small
      Trade time $n$, sign $\varepsilon_n=\pm1$, price in ticks \citep{bouchaud2004fluctuations}:
      \[ P_n = \textstyle\sum_{m<n} \Gamma_0\,(n-m)^{-\zeta}\,\varepsilon_m , \]
      with impact amplitude $\Gamma_0$ (ticks) and decay exponent $\zeta$.
    \end{column}
    \begin{column}{0.48\textwidth}
      \small
      If $C_\varepsilon(\ell)\sim\ell^{-\gamma_\varepsilon}$, then $\mathcal{D}(\ell) = \Var(P_{n+\ell}-P_n)/\ell \sim \ell^{1-\gamma_\varepsilon-2\zeta}$, so prices diffuse only at $\zeta_c = (1-\gamma_\varepsilon)/2$: impact decay offsets sign memory.
    \end{column}
  \end{columns}
  \centering
  \includegraphics[width=0.6\textwidth]{fig6_propagator.png}

  \raggedright\footnotesize
  Signs: sign of fractional Gaussian noise, $\gamma_\varepsilon = \EsixGamma$, $2^{20}$ trades, $\EsixNrep$ replicates. Fitted slopes $\EsixSlopeLow$, $\EsixSlopeCrit$, $\EsixSlopeHigh$ against theory $\EsixTheoryLow$, $\EsixTheoryCrit$, $\EsixTheoryHigh$.
\end{frame}

\begin{frame}{Large orders: the square-root law and the latent book}
  \small
  A metaorder of $Q$ shares, executed over hours, moves the price by
  \[ I(Q) \approx Y\,\sigma_D\,\sqrt{Q/V_D}, \]
  with daily volatility $\sigma_D$ (fraction of price), daily volume $V_D$ (shares) and a dimensionless $Y$ of order one \citep{toth2011anomalous}.
  \begin{itemize}
    \item The visible book cannot explain it: it holds a small fraction of daily volume. \citet{donier2015fully} model a \hl{latent} order book whose density is linear near the price. This gives the square root, contains the propagator model, and allows no price manipulation.
    \item For an LP the message is that impact has two parts: a transient mechanical part that reverts, and a permanent informational part. Markouts at long horizons measure the second.
  \end{itemize}
  \vspace{0.3em}
  The square-root law is an empirical result; this repository does not test it.
\end{frame}
